% --- Archon physics formalization source begin ---
% archon:physics
% archon:covers PhyXMiniProblems/problem_phyx_mini_0039.lean
% archon:source-report reports/phyx_mini/problem_phyx_mini_0039.source.json

\chapter{Physics problem phyx\_mini\_0039}
\label{ch:PhyXMiniProblems_problem_phyx_mini_0039}

\paragraph{Problem source.}
\#\# Physical scenario

Santa checks himself for soot, using his reflection in a silvered Christmas tree ornament 0.750 m away (figure). The diameter of the ornament is 7.20 cm. Standard reference texts state that he is a “right jolly old elf,” so we estimate his height to be 1.6 m.

\#\# Question

How tall is the image of Santa formed by the ornament?

\#\# Answer choices

- A: A: 3.68cm

- B: B: 3.8cm

- C: C: 3.66cm

- D: D: 5.05cm

\#\# Figure caption (auxiliary; use the image as primary evidence)

The image consists of two parts:

1. **Left Side - Illustration of Santa in a Ornament:**
   - The scene is viewed from above, showing a cartoon character dressed as Santa Claus inside a circular ornament.
   - Santa is wearing glasses and his traditional red and white outfit with black boots.
   - He is surrounded by a few gift-wrapped presents on a green floor with a pattern resembling pine branches. 
   - The setting appears to be indoors, possibly near a wall with window-like patterns.

2. **Right Side - Diagram of Optics:**
   - The diagram illustrates a geometric optics setup involving a spherical surface.
   - Labels and text are included:
     - The focal length (f) is given as \textbackslash{}( f = \textbackslash{}frac\{R\}\{2\} = -1.80 \textbackslash{}, \textbackslash{}text\{cm\} \textbackslash{}).
     - The radius of curvature (R) is \textbackslash{}( R = -3.60 \textbackslash{}, \textbackslash{}text\{cm\} \textbackslash{}).
     - The object distance (s) is \textbackslash{}( s = 75.0 \textbackslash{}, \textbackslash{}text\{cm\} \textbackslash{}).
     - The height (y) is \textbackslash{}( y = 1.6 \textbackslash{}, \textbackslash{}text\{m\} \textbackslash{}).
     - The optic axis is marked with a dashed line.
     - The center of curvature is labeled as C.
   - Arrows indicate the direction of light rays interacting with the spherical surface.
   - The annotation "NOT TO SCALE" is included at the bottom.

This combination of illustrations provides a whimsical visual with a physics-based diagram, connecting a playful scene with an optical concept.

\#\# Dataset metadata

Geometrical Optics; Physical Model Grounding Reasoning, Multi-Formula Reasoning

\paragraph{Recorded answer/context.}
B: B: 3.8cm

\paragraph{Figure/image path.}
/root/proposal\_for\_physic/science-mango/phyx\_mini\_run/phyx\_data/test\_image/39.png

\paragraph{Formalization target.}
create a compiling Lean file with sorry bodies at `PhyXMiniProblems/problem\_phyx\_mini\_0039.lean`. The Lean declarations must preserve the physical quantities, dimensions or dimensional roles, figure labels, governing-law hypotheses, and final relation expressed by this problem.
Use Mathlib/Physlib names found through LeanExplore where available. If a physics API is missing, introduce faithful local abstractions rather than scalar placeholder aliases.

\begin{theorem}[Physics formalization target]
\label{thm:physics:phyx_mini_0039:target}
\lean{PhyXMiniProblems.ProblemPhyXMini0039.problem_phyx_mini_0039}
\uses{def:physics:phyx-mini-0039:phyxminiproblems-problemphyxmini0039-lengthincentimeters, def:physics:phyx-mini-0039:phyxminiproblems-problemphyxmini0039-ornamentmirrorsetup, def:physics:phyx-mini-0039:phyxminiproblems-problemphyxmini0039-matchesfigurereadouts, def:physics:phyx-mini-0039:phyxminiproblems-problemphyxmini0039-hasdepictedaxisgeometry, def:physics:phyx-mini-0039:phyxminiproblems-problemphyxmini0039-satisfiesparaxialmirrorlaws, def:physics:phyx-mini-0039:phyxminiproblems-problemphyxmini0039-roundstonearesttenth}
The assigned autoformalize agent should translate this physics problem into Lean declarations in the covered file, with theorem and lemma proof bodies written as `by sorry`.
\end{theorem}
\begin{proof}
This is an autoformalization task, not a proof task. Produce statements that can later be proved by the physics prover without weakening the physical model.
\end{proof}

% --- Archon declaration topology begin ---
\section{Declaration topology}
\begin{definition}[Length Quantity]
  \label{def:physics:phyx-mini-0039:phyxminiproblems-problemphyxmini0039-lengthquantity}
  \lean{PhyXMiniProblems.ProblemPhyXMini0039.LengthQuantity}
  A signed physical length, independent of the unit in which it is read
\end{definition}
\begin{proof}
  This is the declared model definition of Length Quantity.
\end{proof}
\begin{definition}[centimeter Unit Choices]
  \label{def:physics:phyx-mini-0039:phyxminiproblems-problemphyxmini0039-centimeterunitchoices}
  \lean{PhyXMiniProblems.ProblemPhyXMini0039.centimeterUnitChoices}
  Unit choices whose length component is the centimeter used in the figure
\end{definition}
\begin{proof}
  This is the declared model definition of centimeter Unit Choices.
\end{proof}
\begin{definition}[length In Centimeters]
  \label{def:physics:phyx-mini-0039:phyxminiproblems-problemphyxmini0039-lengthincentimeters}
  \lean{PhyXMiniProblems.ProblemPhyXMini0039.lengthInCentimeters}
  \uses{def:physics:phyx-mini-0039:phyxminiproblems-problemphyxmini0039-lengthquantity, def:physics:phyx-mini-0039:phyxminiproblems-problemphyxmini0039-centimeterunitchoices}
  The signed scalar centimeter readout of a dimensionful length
\end{definition}
\begin{proof}
  This is the declared model definition of length In Centimeters.
\end{proof}
\begin{definition}[Spherical Mirror Kind]
  \label{def:physics:phyx-mini-0039:phyxminiproblems-problemphyxmini0039-sphericalmirrorkind}
  \lean{PhyXMiniProblems.ProblemPhyXMini0039.SphericalMirrorKind}
  The two spherical-mirror types as viewed from the incident-light side
\end{definition}
\begin{proof}
  This is the declared model definition of Spherical Mirror Kind.
\end{proof}
\begin{definition}[Axis Location]
  \label{def:physics:phyx-mini-0039:phyxminiproblems-problemphyxmini0039-axislocation}
  \lean{PhyXMiniProblems.ProblemPhyXMini0039.AxisLocation}
  Labeled locations on the common optic axis in the primary figure. imageBasePrime is the base of the virtual image arrow y', and centerOfCurvatureC is the point explicitly labeled C
\end{definition}
\begin{proof}
  This is the declared model definition of Axis Location.
\end{proof}
\begin{definition}[Ornament Mirror Setup]
  \label{def:physics:phyx-mini-0039:phyxminiproblems-problemphyxmini0039-ornamentmirrorsetup}
  \lean{PhyXMiniProblems.ProblemPhyXMini0039.OrnamentMirrorSetup}
  \uses{def:physics:phyx-mini-0039:phyxminiproblems-problemphyxmini0039-lengthquantity, def:physics:phyx-mini-0039:phyxminiproblems-problemphyxmini0039-sphericalmirrorkind, def:physics:phyx-mini-0039:phyxminiproblems-problemphyxmini0039-axislocation}
  The dimensionful quantities in the spherical-mirror diagram. The fields objectDistance and imageDistance are the signed quantities labeled s and s'; objectHeight and imageHeight are the signed transverse quantities labeled y and y'
\end{definition}
\begin{proof}
  This is the declared model definition of Ornament Mirror Setup.
\end{proof}
\begin{definition}[Matches Figure Readouts]
  \label{def:physics:phyx-mini-0039:phyxminiproblems-problemphyxmini0039-matchesfigurereadouts}
  \lean{PhyXMiniProblems.ProblemPhyXMini0039.MatchesFigureReadouts}
  \uses{def:physics:phyx-mini-0039:phyxminiproblems-problemphyxmini0039-lengthincentimeters, def:physics:phyx-mini-0039:phyxminiproblems-problemphyxmini0039-ornamentmirrorsetup}
  Numerical readouts supplied by the problem and the annotations in the primary figure: diameter 7.20 cm, R = -3.60 cm, f = -1.80 cm, object distance s = 75.0 cm, and Santa's height y = 1.6 m = 160 cm. No image distance or image height is specified here
\end{definition}
\begin{proof}
  This is the declared model definition of Matches Figure Readouts.
\end{proof}
\begin{definition}[Has Depicted Axis Geometry]
  \label{def:physics:phyx-mini-0039:phyxminiproblems-problemphyxmini0039-hasdepictedaxisgeometry}
  \lean{PhyXMiniProblems.ProblemPhyXMini0039.HasDepictedAxisGeometry}
  \uses{def:physics:phyx-mini-0039:phyxminiproblems-problemphyxmini0039-lengthincentimeters, def:physics:phyx-mini-0039:phyxminiproblems-problemphyxmini0039-ornamentmirrorsetup}
  Geometry encoded by the figure's optic axis. With coordinates increasing to the right, Santa is in front of the mirror while the virtual image and C are behind it. The signed distances are measured from these labeled positions
\end{definition}
\begin{proof}
  This is the declared model definition of Has Depicted Axis Geometry.
\end{proof}
\begin{definition}[Satisfies Paraxial Mirror Laws]
  \label{def:physics:phyx-mini-0039:phyxminiproblems-problemphyxmini0039-satisfiesparaxialmirrorlaws}
  \lean{PhyXMiniProblems.ProblemPhyXMini0039.SatisfiesParaxialMirrorLaws}
  \uses{def:physics:phyx-mini-0039:phyxminiproblems-problemphyxmini0039-lengthincentimeters, def:physics:phyx-mini-0039:phyxminiproblems-problemphyxmini0039-ornamentmirrorsetup}
  The geometric and paraxial governing laws used for the silvered spherical ornament. They state R = -D/2, f = R/2, the Gaussian mirror equation 1/f = 1/s + 1/s', and the signed transverse-magnification law y'/y = -s'/s. The last two are written without division. These are general modeling relations and contain no numerical image answer
\end{definition}
\begin{proof}
  This is the declared model definition of Satisfies Paraxial Mirror Laws.
\end{proof}
\begin{definition}[Answer Choice]
  \label{def:physics:phyx-mini-0039:phyxminiproblems-problemphyxmini0039-answerchoice}
  \lean{PhyXMiniProblems.ProblemPhyXMini0039.AnswerChoice}
  The four image-height choices printed in the problem
\end{definition}
\begin{proof}
  This is the declared model definition of Answer Choice.
\end{proof}
\begin{definition}[answer Height In Centimeters]
  \label{def:physics:phyx-mini-0039:phyxminiproblems-problemphyxmini0039-answerheightincentimeters}
  \lean{PhyXMiniProblems.ProblemPhyXMini0039.answerHeightInCentimeters}
  \uses{def:physics:phyx-mini-0039:phyxminiproblems-problemphyxmini0039-answerchoice}
  The scalar centimeter value displayed beside each answer choice
\end{definition}
\begin{proof}
  This is the declared model definition of answer Height In Centimeters.
\end{proof}
\begin{definition}[Rounds To Nearest Tenth]
  \label{def:physics:phyx-mini-0039:phyxminiproblems-problemphyxmini0039-roundstonearesttenth}
  \lean{PhyXMiniProblems.ProblemPhyXMini0039.RoundsToNearestTenth}
  \uses{def:physics:phyx-mini-0039:phyxminiproblems-problemphyxmini0039-lengthquantity, def:physics:phyx-mini-0039:phyxminiproblems-problemphyxmini0039-lengthincentimeters, def:physics:phyx-mini-0039:phyxminiproblems-problemphyxmini0039-answerchoice, def:physics:phyx-mini-0039:phyxminiproblems-problemphyxmini0039-answerheightincentimeters}
  The usual half-open bin for rounding a centimeter value to the displayed nearest tenth. Including the lower endpoint makes 3.75 cm round to 3.8 cm, while excluding the upper endpoint keeps the bins disjoint
\end{definition}
\begin{proof}
  This is the declared model definition of Rounds To Nearest Tenth.
\end{proof}
\begin{lemma}[virtual Image Distance eq]
  \label{lem:physics:phyx-mini-0039:phyxminiproblems-problemphyxmini0039-virtualimagedistance-eq}
  \lean{PhyXMiniProblems.ProblemPhyXMini0039.virtualImageDistance_eq}
  \uses{def:physics:phyx-mini-0039:phyxminiproblems-problemphyxmini0039-lengthincentimeters, def:physics:phyx-mini-0039:phyxminiproblems-problemphyxmini0039-ornamentmirrorsetup, def:physics:phyx-mini-0039:phyxminiproblems-problemphyxmini0039-matchesfigurereadouts, def:physics:phyx-mini-0039:phyxminiproblems-problemphyxmini0039-hasdepictedaxisgeometry, def:physics:phyx-mini-0039:phyxminiproblems-problemphyxmini0039-satisfiesparaxialmirrorlaws}
  The Gaussian mirror equation determines the signed virtual-image distance s' = -225/128 cm
\end{lemma}
\begin{proof}
  The conclusion follows from the governing relations and figure readouts named in the statement of virtual Image Distance eq.
\end{proof}
\begin{lemma}[image Height eq]
  \label{lem:physics:phyx-mini-0039:phyxminiproblems-problemphyxmini0039-imageheight-eq}
  \lean{PhyXMiniProblems.ProblemPhyXMini0039.imageHeight_eq}
  \uses{def:physics:phyx-mini-0039:phyxminiproblems-problemphyxmini0039-lengthincentimeters, def:physics:phyx-mini-0039:phyxminiproblems-problemphyxmini0039-ornamentmirrorsetup, def:physics:phyx-mini-0039:phyxminiproblems-problemphyxmini0039-matchesfigurereadouts, def:physics:phyx-mini-0039:phyxminiproblems-problemphyxmini0039-hasdepictedaxisgeometry, def:physics:phyx-mini-0039:phyxminiproblems-problemphyxmini0039-satisfiesparaxialmirrorlaws}
  The mirror equation and transverse-magnification law determine the exact upright image height y' = 15/4 cm
\end{lemma}
\begin{proof}
  The conclusion follows from the governing relations and figure readouts named in the statement of image Height eq.
\end{proof}
% --- Archon declaration topology end ---
% --- Archon physics formalization source end ---
